\section{Comparison with the double-cover count on a closed threefold}\label{sec:normalization}
We turn to a second compact consequence. The chamber weight for smooth double covers is normalized at a complex-linear operator. In the elementary cluster expansion the reference is the balanced local curve. Computing the two references gives a correction in twice the primitive class.
Let $v$ be the genus parameter. Write $s=2\sin(v/2)$ and $x=s^2$. For a partition $\lambda$ and a box
$u$ of its diagram, let $\ell(u)$ denote its hook length. The connected
elementary series of a curve of genus $g$ is
\[
 G_g(q,v)=\log\!\left(\sum_{\lambda}q^{|\lambda|}
            \prod_{u\in\lambda}(2\sin(\ell(u)v/2))^{2g-2}\right).
\]
The empty diagram contributes one before taking the logarithm.
Ionel--Parker's $e^{\rm IP}_{B,h}$ is the coefficient of
$G_h(Q^B,v)$ in the connected Gromov--Witten series
\cite[Proposition~3.3 and (8.1)]{ip}. In particular
$[q]G_g=x^{g-1}$. The normal weights in the elementary series are
opposite. The class variable $Q$ records integral homology classes.

For a generic compatible almost complex structure, denote the
determinant sign of the normal operator $D_C$ of an embedded curve $C$
by $\epsilon_C$. The Bai--Swaminathan count in genus $h$ is
\[
 \operatorname{Gr}^{\rm BS}_{2A,h}
 =\sum_{\substack{[C]=2A\\g(C)=h}}\epsilon_C
       +\sum_{[C]=A}\epsilon_C w_{2,h}(D_C).
\]
Here $w_{2,h}$ is the oriented intersection number with the wall of
smooth genus-$h$ double covers carrying an anti-invariant kernel
\cite[(1.5), Definitions~5.8 and~5.17]{bai}. Its value at a
double-cover-rigid complex-linear reference is zero. We use the
coorientation and path convention of that definition throughout.
We extend $w_{2,h}$ by zero when $h<g(C)$.

\begin{theorem}\label{p9-thm:global}
Let $X$ be a closed symplectic six-manifold with $c_1(TX)=0$ and torsion-free $H_2(X;\ZZ)$. For every primitive class $A$ of
positive area and every $h\ge0$,
\begin{equation}\label{p9-eq:global}
 e^{\rm IP}_{2A,h}(X)-\operatorname{Gr}^{\rm BS}_{2A,h}(X)
       =\sum_{g\ge0} e^{\rm IP}_{A,g}(X)\,a_{g,h}.
\end{equation}
The coefficient $a_{g,h}$ vanishes for $g=0,1$. For $c=g-1\ge1$ and
$j=h-(2g-1)$ it is
\begin{equation}\label{p9-eq:coeff}
 a_{g,h}=
 \begin{cases}
  2\,4^{c-j}\left(\binom{c}{2j}-(-1)^j\binom{c}{j}\right),&1\le j\le c,\\[2pt]
  0,&\text{otherwise}.
 \end{cases}
\end{equation}
A binomial coefficient outside its usual range is zero.
\end{theorem}

The first two nonzero rows of the correction, written in the variable
$x$ with genus $h$ contributing $x^{h-1}$, are $2x^3$ and
$24x^5-2x^6$. On the compact threefold of Section~\ref{p9-sec:example},
a primitive class $A$ has $e^{\rm IP}_{A,3}=-14$ and no higher-genus
primitive coefficient. Only the term $(-14)(-2)$ contributes in
genus seven. Corollary~\ref{p9-cor:example} gives a family on the same
threefold for which the absolute differences grow quadratically.

\subsection{The fixed-target comparison}
Let $C$ be a smooth closed connected Riemann surface of genus $g$ and
let $N\to C$ be a complex bundle of rank two and degree $2g-2$. Let $D$
be a real Cauchy--Riemann operator on $N$. Its index on $C$ is zero.
When $D$ is invertible, $\epsilon(D)$ is its sign for the determinant
orientation obtained by transport from complex-linear operators.
Suppose also that $f^*D$ is injective for every smooth connected
double cover $f$ of $C$. We call such an operator rigid through degree
two. On a nodal stable map of degree two one can normalize the positive
components, each of degree one or two. Kernel sections vanish there,
and the constants on each connected contracted subcurve vanish by
matching. Hence all stable pullbacks are injective.

Write $M_k=\overline{\mathcal M}_k(C,2)$ for the stack of connected
degree-two stable maps with unmarked source of arithmetic genus $k$.
The oriented real cokernels of $f^*D$ give a bundle on $M_k$ of real rank
$2r$, where $r=2k-2-2(2g-2)$ is the complex virtual dimension of $M_k$. Define
\[
 L_{2,k}(D)=\int_{[M_k]^{\rm vir}} e(\operatorname{coker} f^*D),
 \qquad B_{2,g}^{\rm bal}(v)=[q^2]G_g(q,v).
\]
The superscript $\rm bal$ denotes the balanced antidiagonal series.
The integral is zero for empty moduli. We use the orientation of the
full pullback determinant for $L_{2,k}$, including its sign in rank
zero. In the orientation relative to the base operator, this number
is $\epsilon(D)L_{2,k}(D)$ \cite[Remark~6.4]{bai}.

\begin{theorem}\label{p9-thm:local}
For every operator rigid through degree two one has
\begin{equation}\label{p9-eq:local}
 \epsilon(D)\sum_{k\ge0}L_{2,k}(D)v^{2k-2}
    = B_{2,g}^{\rm bal}(v)+
       \sum_{h\ge0}\bigl(w_{2,h}(D)+a_{g,h}\bigr)x^{h-1}.
\end{equation}
The equality is coefficientwise in $v$. In genus zero the series may
have a pole of order two.
\end{theorem}

\subsection{Comparison with the chamber weight}\label{p9-sec:comparison}

\subsubsection{Transport from a complex-linear operator}
Fix $C$ and $N$ as in Theorem~\ref{p9-thm:local}. Use an invertible
complex-linear operator $\beta$ rigid through degree two as a reference.
Such operators exist by the complex-operator transversality in
\cite[Lemmas~5.3, 5.6 and~5.15(ii), Definition~5.17]{bai}. Base acyclicity is
open and is achieved by taking a direct sum of degree-$g-1$ line
bundles off the theta divisor, or $\mathcal O(-1)^{\oplus2}$ when
$g=0$. The reference need not split holomorphically.

Choose a generic path from $\beta$ to $D$ in the affine space of real
Cauchy--Riemann operators with fixed symbol. At an ordinary crossing
of the base wall the kernel is a real line of nowhere-zero sections.
At an ordinary crossing of a double-cover wall its anti-invariant
kernel line gives an embedded graph and its intrinsic derivative
\eqref{p9-eq:wall-differential} is invertible. Here one uses the
one-parameter transversality and embeddedness of the kernel graphs
\cite[Lemmas~5.6, 5.11, 5.15--5.16]{bai}. The variations in these tests
are variations of the zeroth-order real tensor on an unramified disk.
The fixed-symbol space also has the proper degree-at-most-two wall of
\cite[Definition~2.6, Proposition~2.7 and Theorem~3.40]{dw}. In
particular, up to any prescribed genus there are only finitely many
crossings.

We give the transversality detail that separates a base crossing from
an odd crossing. Mark real kernel lines on the base and in the
anti-invariant part of a smooth cover with $b$ branch points. For a
fixed cover the Fredholm indices of the two operators are $0$ and
$-2b$. Projectivizing both kernels and allowing the branch values to
vary gives index $-2$. For surjectivity of the universal section,
an adjoint pair annihilating all variations of the operators would,
on an unramified disk, annihilate all pairs of values
$(A\xi_0,A\xi_1)$, where $A$ is a localized real tensor and the two
generators are evaluated on one sheet. The odd kernel and its
adjoint both change sign on the second sheet, so the contributions
of the sheets add. The generators are real independent on a disk.
Otherwise on a disk of nonvanishing a real multiple relating two
solutions has zero $\bar\partial$, hence is constant. Unique
continuation on the connected cover then contradicts their opposite
deck characters. Arbitrary variations of $A$ on the disk force both
adjoints to vanish there and hence globally. Thus a generic path
misses this incidence. Taking the countable intersection over source
genera preserves genericity.

Put $E_{2,k}(D)=\epsilon(D)L_{2,k}(D)$ on regular slices. At a base
crossing with no odd kernel this reduced Euler number extends across
the crossing. To see this, for a smooth positive component of degree
two the invariant sections in the pullback kernel descend away from the
branch values. They are bounded there in a target trivialization.
Removable singularities for the real Cauchy--Riemann equation extend
them across the branch values. Since there is no odd kernel, all kernel
sections are therefore pulled back from the base. If there are two
degree-one components, their base sections agree at a point, either at
a common node or along the chain of contracted components joining
them. Their difference vanishes everywhere because a nonzero base
kernel section is nowhere zero. Constants on each contracted subcurve
are determined by their nodal values. The same description on a single
positive component is compatible with any further node. Consequently
the kernel on every stable map is precisely the base pullback.

Represent the base and pullback families by compatible stabilized
two-term complexes of finite-dimensional bundles as in
\cite[Appendix~B, (B.7)--(B.10)]{bai}. The reduced complex is the
quotient in this finite-dimensional presentation. To see that its
degree-zero cohomology vanishes, it remains, by the cohomology exact
sequence and the kernel computation above, to check that pullback
is injective on cokernels. Represent a base cokernel class by a smooth
form $\alpha$. If $f^*D u=f^*\alpha$ on a single effective double
normalization, averaging $u$ over the two sheets gives a base solution
of $D\bar u=\alpha$ off the branch values. It extends by bounded
removability for the inhomogeneous equation. On two degree-one components
either restriction supplies such a solution. Thus the base class
vanishes. The finite-dimensional reduced differential is injective
over the whole compact stable-map space at the crossing.
Its cokernels form an oriented bundle across the base crossing. The
determinant identity is
\[
 \det(f^*D)\simeq\det D\otimes\det((f^*D)_{\rm red}).
\]
In particular the reduced integral on either side is $E_{2,k}$.

At a double-cover crossing the base sign is constant. In the convention
of \cite[(5.59)--(5.60)]{bai}, the forward change of $w_{2,H}$ is
$\epsilon(D_0)\eta_f$ for a cover of genus $H$. Applying
Theorem~\ref{p9-thm:wall} and summing the jumps within each source genus
gives
\begin{equation}\label{p9-eq:transport}
 \sum_{k\ge0}\bigl(E_{2,k}(D)-E_{2,k}(\beta)\bigr)v^{2k-2}
       =\sum_{h\ge0} w_{2,h}(D)x^{h-1}.
\end{equation}
At the initial endpoint $w_{2,h}(\beta)=0$ by the constant path in its
definition. The equality follows for every formal coefficient using
compactness only through that coefficient's genus.

\subsubsection{The Euler integral at the reference}
Let $\pi$ and $f$ now denote the universal curve and evaluation map
over $M_k$. At $\beta$ the class $-R\pi_* f^*N$ is represented by its
actual complex cokernel. Its rank is $r=\dim_{\CC}^{\rm vir}M_k$ and
the base determinant has positive sign. Thus its ordinary top Chern
number is $E_{2,k}(\beta)$. To compute it give $N$ scalar torus weight
$t$. If $V$ is the virtual index class of rank $r$, its equivariant Euler
class is $t^r c_{1/t}(V)$, with $c_z$ the total Chern polynomial.
Integration on the virtual cycle selects $c_r(V)$. This description
also permits deformation of the holomorphic structure to a sum of
two lines of degree $g-1$, since the family index class extends along
the deformation, including at split operators with kernel.

Equal weights compute the reference number, whereas opposite weights
compute the elementary contribution already assigned to the base curve.
The scalar and antidiagonal specializations have the same degree-one
term $x^{g-1}$. In degree two taking the connected logarithm subtracts
the same half-square from both. The scalar formula of
\cite[\S8]{bp08} and the two antidiagonal hook products in
\cite[Corollary~7.2]{bp08} therefore give
\begin{align*}
 \sum_{k\ge0} E_{2,k}(\beta)v^{2k-2}-B_{2,g}^{\rm bal}(v)
  =x^{2g-2}\bigl((4-2s)^{g-1}+(4+2s)^{g-1}
                           -2(4-x)^{g-1}\bigr).
\end{align*}
The bracket vanishes for $g=1$. For $g=0$ the sum of the first
two reciprocals is $2/(4-x)$. For $c=g-1>0$ the odd powers of $s$
cancel. The coefficient of $x^{2c+j}$ is \eqref{p9-eq:coeff} for
$h=2g-1+j$, and the constant term in the bracket vanishes. Together
with \eqref{p9-eq:transport} this proves Theorem~\ref{p9-thm:local}.

\subsubsection{Passage to the compact invariants}
Choose a generic compatible almost complex structure for which all
simple curves of index zero are embedded, regular, super-rigid and
pairwise disjoint \cite[\S1]{bai}. The image of a connected stable
map in class $2A$ lies on a single such support. Indeed, two distinct
effective images cannot meet, and contracted components can only join
them at a common image point. If the map has degree $m$ on its support
of class $B$, then $mB=2A$. Primitivity and the torsion-free lattice
give $m=1$, $B=2A$, or $m=2$, $B=A$. The same argument in class $A$
gives only degree one.

Through any fixed source genus there are finitely many contributing
simple supports, by compactness at fixed genus and super-rigidity
\cite[\S\S3,5]{bai}. The degree-one series on a support of genus $g$
is $\epsilon_C x^{g-1}$, by the degree-one calculation in
Theorem~\ref{p9-thm:wall}. Hence $e^{\rm IP}_{A,g}$ is the sum of signs
of genus-$g$ supports of class $A$. In class $2A$ first subtract their
degree-two elementary terms $\epsilon_C B_{2,g}^{\rm bal}$. By
\eqref{p9-eq:local}, the remaining coefficient of $x^{h-1}$ on each such
support is $\epsilon_C(w_{2,h}(D_C)+a_{g,h})$. A support in $2A$
contributes its degree-one determinant sign in its own genus. The
latter signs and the terms $\epsilon_C w_{2,h}$ give exactly
\cite[(1.5)]{bai}. The remaining terms sum to \eqref{p9-eq:global},
proving Theorem~\ref{p9-thm:global}.
\subsection{A compact threefold with unbounded discrepancy}\label{p9-sec:example}

Let $X_9$ be a general nondegenerate anticanonical hypersurface in the smooth
projective toric fourfold with rays and maximal cones
\cite[(23)--(24)]{WueCounter}. We equip it with a K\"ahler form.
It has trivial canonical bundle and torsion-free $H_2(X_9;\ZZ)$
\cite[\S4.1]{WueCounter}. The divisor corresponding to ray $8$ cuts out
a surface $T\simeq\mathbb F_1$. Write $s$ for its negative section, $f$ for a
ruling fiber and $i:T\hookrightarrow X_9$ for the inclusion. For $d\ge3$ let
$D_d=ds+(d+1)f$ and $A_d=i_*D_d$. Set
\begin{equation}\label{p9-eq:compact-classes}
 G=\frac{d(d-1)}2,\qquad N=\frac{(d+1)(d+4)}2 .
\end{equation}

\begin{corollary}\label{p9-cor:example}
The class $A_d$ is primitive, and
\[
 e^{\rm IP}_{2A_d,3G-2}(X_9)-\operatorname{Gr}^{\rm BS}_{2A_d,3G-2}(X_9)
     =(-1)^{d+1}(d+1)(d+4).
\]
In particular, the two counts differ by $28$ in class $2i_*(3s+4f)$
and genus seven. Their differences are unbounded on $X_9$.
\end{corollary}

\subsubsection{The surface classes}
The semiample divisor $H_4$ of \cite[\S4.2]{WueCounter} contracts curves
only in $T$, in a second surface $E\simeq\operatorname{Bl}_{p_1,p_2}\PP^2$,
or in two rational curves $C_1,C_2$. On $E$ the effective cone is generated
by its exceptional curves $e_1,e_2$ and the strict transform of the line
through the two points, which equals $s=E\cap T$. The six integral nef
divisors of \cite[Lemma~4.1 and (30)]{WueCounter} have degrees
\[
\begin{array}{c|rrrrrr}
 & H_1&H_2&H_3&H_4&H_5&H_6\\ \hline
 s   &0&1&0&0&0&0\\
 f   &0&0&1&0&0&0\\
 e_1 &1&0&0&0&1&1\\
 e_2 &0&0&0&0&0&1\\
 C_1 &0&0&0&0&1&0\\
 C_2 &1&0&0&0&0&0
\end{array}
 .
\]
In particular $H_3-H_2$ pairs to one with $A_d$, and $H_2,H_3$
distinguish all classes on $T$. Each effective summand of $A_d$ has zero
$H_1,H_4,H_5,H_6$ degrees. These conditions exclude $C_1,C_2$ and force
any component on $E$ to have class proportional to its negative curve $s$.
Such a component equals $s$. Thus every summand is represented entirely
on $T$. This is the argument of \cite[Proposition~4.2]{WueCounter}
for arbitrary $d$.

\subsubsection{The primitive coefficient}
Adjunction gives $N_{T/X_9}=K_T$. To compute the connected invariants on
these classes, we deform $T\subset X_9$ to $Z=\operatorname{Tot}(K_T)$. Blowing up $T\times 0$ in $X_9\times\mathbb A^1$ and removing
the strict transform of the central fiber gives a smooth family $V$ with
fibers $X_9$ away from zero and $Z$ over zero. It contains the Cartier
divisor $\mathcal T=T\times\mathbb A^1$ and has a natural morphism to $X_9$.
Use this morphism to label curve classes. In $Z$ every complete irreducible
curve lies in $T$. Its projection must be nonconstant since the fibers are
affine lines, and the tautological fiber coordinate on its normalization
is a section of a line bundle of negative degree.

For every positive class on this face, stable maps from connected domains to
the fibers of $V$ factor through $\mathcal T$ in families. Indeed on a
geometric fiber the image lies in $T$ by the preceding support calculation.
The pullback of $\mathcal O_V(\mathcal T)$ then has negative degree on
each nonconstant component and is trivial on each contracted component.
Matching at the nodes makes its space of sections zero. For a family of
nodal domains the derived direct image of this line bundle is locally a
two-term complex of vector bundles in degrees zero and one. Fiberwise
injectivity of its differential implies that its degree-zero cohomology
sheaf vanishes after every base change. The defining section of $\mathcal T$
therefore pulls back to zero, giving the asserted factorization also over
nonreduced bases. The relative stable-map stack in the specified class
and genus is consequently the product of a proper stable-map stack of $T$
with $\mathbb A^1$. Deformation invariance of its relative virtual class
identifies the connected invariants of $X_9$ and $Z$ in every genus.

Write $D=D_d$. On $T=\operatorname{Bl}_p\PP^2$ it is the class of a plane
curve of degree $d+1$ through $p$ with one imposed simple base point.
Thus $p_a(D)=G$ and $h^0(T,\mathcal O_T(D))=N$. We claim that the least
power of $z$ in the stable-pair partition function of $Z$ in class $D$ is
$z^{1-G}$, with coefficient $(-1)^{N-1}N$, and that this term survives
in the logarithm. The case $d=3$ of the following moduli calculation appears
in \cite[Propositions~5.2--5.3]{WueCounter}.

For any splitting $D=\sum_j D_j$ into at least two nonzero effective classes
on $T$, write $D_j=a_js+b_jf$ with $a_j,b_j\ge0$. From
$s^2=-1$, $sf=1$ and $f^2=0$ we get
\[
 \sum_j\chi(\mathcal O_{D_j})-\chi(\mathcal O_D)
  =\sum_{j<k}D_jD_k
  =\sum_j b_j(d-a_j)-\sum_{j<k}a_ja_k>0.
\]
Indeed, let $m=\max a_j$ and $r=d-m$. If $r>0$, the positive sum is
at least $(d+1)r$, and the sum being subtracted is at most
$mr+r(r-1)/2$. The difference is at least $r(r+3)/2$.
If $r=0$, the remaining terms are positive multiples of $f$, giving a
difference of at least $d$.

The image of a stable pair is the structure sheaf of a pure curve $B$,
and its cokernel has finite length. Filter $\mathcal O_B$ by powers of
the ideal of $T$. Its $j$th quotient is
$K_T^{-j}\otimes(\mathcal O_T/J_j)$ for some ideal $J_j$, since the
conormal bundle is $K_T^{-1}$. Write
$J_j^{\vee\vee}=\mathcal O_T(-D_j)$ and
$\ell_j=\operatorname{length}(J_j^{\vee\vee}/J_j)$. Then
\[
 \chi(\mathcal O_B)=
 \sum_j\bigl(\chi(\mathcal O_{D_j})-jK_TD_j+\ell_j\bigr),
 \qquad \sum_j[D_j]=[B].
\]
The correction terms are nonnegative because $-K_T$ is ample. If there
are several divisorial layers, the strict inequality above gives
$\chi(\mathcal O_B)>1-G$ in class $D$. If there is only one, it is the
zeroth layer by Nakayama's lemma at the generic points. The ideal of
$T$ then generates a subsheaf of $\mathcal O_B$ of finite support, which
vanishes by purity. Hence equality occurs precisely for quotient pairs
$\mathcal O_Z\to\mathcal O_B$ with $B\in|D|$. For a product of pairs
in two or more nonzero classes summing to $D$, their combined filtrations
have at least two nonzero divisorial layers. The same strict inequality
shows that all product terms in the logarithm occur above $z^{1-G}$.

The quotient pairs give a closed $\PP^{N-1}=|D|$ in the minimum
pair space. They exhaust its closed points, and surjectivity of their
sections persists in every infinitesimal deformation by Nakayama's lemma.
For every $B\in|D|$, including nonreduced $B$, the normal sequence is
\[
 0\longrightarrow\mathcal O_B(D)\longrightarrow N_{B/Z}
   \longrightarrow K_T|_B\longrightarrow0.
\]
Now $H^0(K_T)=0$ and $H^1(K_T-D)=H^1(\mathcal O_T(D))^*=0$ by Kodaira
vanishing, since $D-K_T=(d+2)s+(d+4)f$ is ample.
Thus $H^0(K_T|_B)=0$, while $h^0(\mathcal O_B(D))=N-1$ by rationality
of $T$. The Hilbert tangent dimension is $N-1$ at every closed point.
It equals the local dimension, so the pair space is regular and the closed
immersion is an isomorphism. Its symmetric obstruction theory has virtual
degree $(-1)^{N-1}\chi_{\rm top}(\PP^{N-1})=(-1)^{N-1}N$
\cite[\S1]{pt}.

Finally, apply the toric Gromov--Witten/stable-pairs correspondence
\cite[Theorem~1]{pp} to $Z$, with empty insertions and $-z=e^{iu}$.
The Calabi--Yau condition gives prefactor one. The stable-map and
stable-pair spaces in these classes are proper. For pairs, compactify $Z$
to the projective line bundle over $T$. In a flat family over a valuation
ring the support is the closure of the generic support, so a pair
supported set-theoretically on the zero section specializes there as well.
Thus the open space of pairs supported in $Z$ is proper, and the
equivariant integrals give ordinary virtual degrees. The disconnected
Gromov--Witten series in this correspondence omits degree-zero connected
components, so taking its logarithm gives the connected series.
The equality applies to every effective summand on the face.

On the projective fiber, the algebraic and symplectic virtual degrees
agree \cite{siebertgw}. By primitivity, the coefficient of class $A_d$ in
the connected elementary expansion is the finite sum of degree-one terms
\cite[Proposition~3.3 and Theorem~8.1]{ip}, \cite{diw}. It follows that
\[
 [Q^D]\log Z_{\rm PT}^{Z}(z,Q)
     =\sum_{g\ge0} e^{\rm IP}_{A_d,g}(X_9)
           \left(\frac{(1+z)^2}{z}\right)^{g-1}.
\]
Both sides are rational in $z$. The correspondence identifies their expansions
at $-z=e^{iu}$, since $(2\sin(u/2))^2=(1+z)^2/z$, and therefore also
their Laurent expansions at zero. Each summand of genus $g$ has lowest
power $z^{1-g}$ with leading factor one. The minimum computed above gives
\begin{equation}\label{p9-eq:compact-primitive}
 e^{\rm IP}_{A_d,G}(X_9)=(-1)^{N-1}N,\qquad
 e^{\rm IP}_{A_d,g}(X_9)=0\quad(g>G).
\end{equation}

\begin{proof}[Proof of Corollary~\ref{p9-cor:example}]
In \eqref{p9-eq:global} a row with base genus $g<G$ vanishes in
genus $3G-2$ by \eqref{p9-eq:coeff}. All rows with $g>G$ vanish by
\eqref{p9-eq:compact-primitive}. The remaining coefficient in
\eqref{p9-eq:coeff} is $a_{G,3G-2}=-2(-1)^{G-1}$. Multiplication by
$(-1)^{N-1}N$ gives the formula, since $G+N-2=d(d+2)$.
For $d=3$ the minimum pair space is $\PP^{13}$ and its virtual degree
is $-14$, giving $(-2)(-14)=28$. For $d=15$ the difference is $304$
in genus $313$. As $d$ ranges through $12\ell+3$ these positive differences
tend to infinity.
\end{proof}
